\documentclass[10pt,aspectratio=169]{beamer}
\usetheme[progressstyle=movingCircCnt]{Feather}

\usepackage[utf8]{inputenc}
\usepackage[T1]{fontenc}
\usepackage[english,brazil]{babel}
\usepackage{helvet}
\usepackage{amsmath}
\usepackage{booktabs}
\usepackage{tikz}
\usetikzlibrary{arrows.meta,positioning}

\definecolor{unbblue}{RGB}{0,51,102}
\definecolor{unbgreen}{RGB}{33,140,33}
\definecolor{softgreen}{RGB}{230,244,230}
\setbeamercolor{alerted text}{fg=unbgreen}
\setbeamerfont{frametitle}{size=\large,series=\bfseries}
\setbeamertemplate{itemize item}{\color{unbgreen}$\blacktriangleright$}
\setbeamertemplate{itemize subitem}{\color{unbgreen}$\bullet$}

\title[Cifra de Vigenère]{\textbf{Cifra de Vigenère:\\Implementação e Criptoanálise}}
\subtitle[CIC0201]{Trabalho de Implementação 1 -- CIC0201}
\author[Grupo Vigenère]{Ana Luísa Reis Nascente \and Gabriel de Sousa\\Marina Pimentel Moreno \and Pedro de Paula Campos}
\institute[]{Segurança Computacional\\Universidade de Brasília}
\date{1 de setembro de 2026}

\newcommand{\source}[1]{\vfill{\tiny\color{gray}Fonte: #1}}

\begin{document}

{\1
\begin{frame}[plain,noframenumbering]
  \titlepage
\end{frame}}

\section{Introdução e contexto}
\begin{frame}{Introdução e contexto}{Por que estudar Vigenère?}
  \begin{columns}[T,onlytextwidth]
    \column{0.57\textwidth}
      \begin{itemize}
        \item Cifra clássica de \textbf{substituição polialfabética}: cada posição pode usar um deslocamento diferente.
        \item Uma chave de comprimento $L$ cria um espaço de $26^L$ possibilidades.
        \item Porém, a \textbf{repetição cíclica} da chave mantém padrões estatísticos da linguagem.
      \end{itemize}
    \column{0.39\textwidth}
      \begin{block}{Objetivo do trabalho}
        Implementar a cifra do zero e recuperar \textbf{tamanho da chave}, \textbf{chave} e \textbf{texto claro} sem usar bibliotecas criptográficas prontas.
      \end{block}
  \end{columns}
  \source{Roteiro do Trabalho de Implementação 1 -- CIC0201, UnB (2026).}
\end{frame}

\section{Técnica teórica}
\begin{frame}{Técnica teórica}{Cifração e decifração}
  \begin{columns}[T,onlytextwidth]
    \column{0.50\textwidth}
      \begin{block}{Vigenère}
        \centering
        $C_i=(P_i+K_{i\bmod L})\bmod 26$\\[8pt]
        $P_i=(C_i-K_{i\bmod L}+26)\bmod 26$
      \end{block}
      \vspace{4pt}
      A chave é repetida sobre as letras do texto, produzindo uma sequência de cifras de César.
    \column{0.46\textwidth}
      \begin{block}{Política adotada}
        \begin{itemize}
          \item somente A--Z e a--z são cifrados;
          \item maiúsculas e minúsculas são preservadas;
          \item espaços, números, acentos e pontuação não mudam;
          \item esses caracteres não avançam a chave.
        \end{itemize}
      \end{block}
  \end{columns}
\end{frame}

\begin{frame}{Técnica teórica}{Como o ataque recupera a chave}
  \begin{columns}[T,onlytextwidth]
    \column{0.48\textwidth}
      \begin{block}{1. Índice de Coincidência}
        \centering
        $IC=\dfrac{\sum_i n_i(n_i-1)}{N(N-1)}$
      \end{block}
      \begin{itemize}
        \item linguagem natural: aproximadamente 0,065--0,075;
        \item distribuição uniforme: aproximadamente $1/26=0,038$;
        \item o maior IC médio indica períodos prováveis da chave.
      \end{itemize}
    \column{0.48\textwidth}
      \begin{block}{2. Análise de frequência}
        \begin{itemize}
          \item divide o criptograma por posição da chave;
          \item cada grupo vira uma cifra de César;
          \item testa os 26 deslocamentos;
          \item compara com frequências de português ou inglês e reconstrói a chave.
        \end{itemize}
      \end{block}
  \end{columns}
  \source{Friedman, \textit{The Index of Coincidence and Its Applications in Cryptography} (1922).}
\end{frame}

\begin{frame}{Técnica teórica}{Pipeline integrado}
  \centering
  \begin{tikzpicture}[
    node distance=5mm,
    box/.style={draw=unbblue, rounded corners, fill=softgreen, minimum height=12mm,
      text width=22mm, align=center, font=\small},
    arrow/.style={-{Stealth[length=2.5mm]}, thick, draw=unbgreen}
  ]
    \node[box] (in) {Criptograma};
    \node[box,right=of in] (ic) {Estimar $L$\\pelo IC};
    \node[box,right=of ic] (freq) {Analisar\\frequências};
    \node[box,right=of freq] (key) {Reconstruir\\a chave};
    \node[box,right=of key] (out) {Decifrar e\\validar};
    \draw[arrow] (in) -- (ic);
    \draw[arrow] (ic) -- (freq);
    \draw[arrow] (freq) -- (key);
    \draw[arrow] (key) -- (out);
  \end{tikzpicture}

  \vspace{12pt}
  \begin{columns}[T,onlytextwidth]
    \column{0.48\textwidth}
      \textbf{Componentes:} cifrador/decifrador, estimador, frequências PT/EN e ataque automático.
    \column{0.48\textwidth}
      \textbf{Saída:} vários candidatos ordenados para avaliar legibilidade e períodos equivalentes.
  \end{columns}
\end{frame}

\section{Resultados}
\begin{frame}{Resultados}{Estratégia de validação}
  \begin{columns}[T,onlytextwidth]
    \column{0.55\textwidth}
      \begin{itemize}
        \item Projeto unificado em C++17 e CMake.
        \item Testes do cifrador, CLI, preservação de caracteres e entradas inválidas.
        \item 500 casos determinísticos de cifrar e decifrar.
        \item Testes do estimador com chaves de tamanhos 5, 7, 11 e 16.
        \item Teste integrado com texto conhecido em português.
      \end{itemize}
    \column{0.39\textwidth}
      \begin{block}{Resultado da suíte}
        \centering
        {\fontsize{28}{32}\selectfont\color{unbgreen}\textbf{7/7}}\\[4pt]
        testes aprovados
      \end{block}
      \begin{block}{Entrada do ataque}
        Idioma: PT\\
        Limite: 20\\
        Chave desconhecida pelo pipeline
      \end{block}
  \end{columns}
\end{frame}

\begin{frame}{Resultados}{Experimento controlado em português}
  \begin{columns}[T,onlytextwidth]
    \column{0.52\textwidth}
      \begin{block}{Primeiro candidato}
        \centering
        \textbf{Tamanho: 9}\\[5pt]
        {\Large\ttfamily\color{unbgreen}SEGURANCA}
      \end{block}
      \begin{itemize}
        \item texto claro recuperado integralmente;
        \item espaços, pontuação e parágrafos preservados;
        \item resultado igual ao arquivo de referência.
      \end{itemize}
    \column{0.44\textwidth}
      \begin{alertblock}{Período equivalente}
        O segundo candidato foi tamanho 18 e chave
        \begin{center}
          \ttfamily SEGURANCASEGURANCA
        \end{center}
        Ele também decifra corretamente por repetir duas vezes a chave fundamental.
      \end{alertblock}
      \textbf{Decisão:} priorizar a menor chave equivalente.
  \end{columns}
  \source{Arquivos e execução registrados em \texttt{textos para testes/README.md}.}
\end{frame}

\section{Conclusões}
\begin{frame}{Conclusões}{O que o experimento demonstrou}
  \begin{itemize}
    \item A implementação cifra e decifra corretamente, preservando a política de caracteres definida.
    \item O pipeline recuperou \textbf{tamanho 9}, \textbf{chave SEGURANCA} e o \textbf{texto claro completo}.
    \item Um espaço de chaves grande não garante segurança quando uma chave curta é reutilizada ciclicamente.
    \item O desempenho da análise depende do tamanho do texto, do idioma correto e da escolha entre períodos equivalentes.
  \end{itemize}

  \vspace{6pt}
  \begin{block}{Conclusão central}
    A vulnerabilidade explorada está na \textbf{repetição da chave}: uma chave aleatória, do tamanho da mensagem e usada uma única vez (One-Time Pad) não apresenta esse mesmo padrão estatístico.
  \end{block}
\end{frame}

\end{document}
